\section{Related Work}
\label{sec:related_work}

\noindent\textbf{Traditional Feature-Based Approaches.}
Early automated DaT-SPECT classification relied on region-of-interest (ROI) analysis and specific binding ratios (SBR). The caudate and putamen are segmented using anatomical atlases (e.g., Hammers, AAL, or probabilistic striatal atlases), and uptake ratios relative to a reference region (occipital cortex or cerebellum) are computed. Support vector machines (SVM), random forests, and gradient boosting machines (GBM) trained on these features achieved AUCs of 0.85--0.92~\cite{segovia2019dat,liu2020dat,khan2021dat}. While interpretable, these methods depend critically on accurate registration and atlas quality, and discard spatial patterns outside predefined ROIs.

\noindent\textbf{Deep Learning for 3D Medical Image Classification.}
The success of 2D CNNs in natural images motivated 3D extensions for volumetric medical data. Key architectures include VoxNet~\cite{maturana2015voxnet}, 3D ResNet~\cite{hara20183d}, DenseNet~\cite{li2018densely}, and EfficientNet3D~\cite{chen2020efficientnet3d}. For DaT-SPECT, Choi et al.~\cite{choi2018deep} trained a 3D CNN on 91$^3$ volumes achieving 0.93 AUC. Kim et al.~\cite{kim2020deep} used a 3D DenseNet with attention mechanisms reaching 0.95 AUC on a single-site dataset. However, most studies evaluate on homogeneous single-site data, limiting generalizability.

\noindent\textbf{Multi-Site and Domain Adaptation.}
Multi-site heterogeneity remains a major challenge. Domain adaptation techniques (adversarial training~\cite{ganin2016domain}, style transfer~\cite{zhang2020style}, batch normalization adaptation~\cite{li2019adaptive}) have been explored for MRI and CT, but less so for SPECT. Ensemble methods combining models trained on different sites or resolutions show promise~\cite{li2020ensemble}. Our stratified group cross-validation by scanner site provides a rigorous alternative that directly measures generalization to unseen sites without requiring adaptation.

\noindent\textbf{Calibration and Uncertainty.}
Modern neural networks are often overconfident~\cite{guo2017calibration}. Temperature scaling~\cite{guo2017calibration}, Platt scaling~\cite{platt1999probabilistic}, and isotonic regression~\cite{zadrozny2002transforming} are post-hoc calibration methods. For medical applications, calibration is essential for risk stratification and clinical decision-making~\cite{kull2019beyond}. We employ Platt scaling on OOF ensemble predictions, achieving significant LogLoss improvement (0.2618 $\to$ 0.2453) while preserving AUROC.

\noindent\textbf{Test-Time Augmentation (TTA).}
TTA improves robustness by averaging predictions over transformed inputs. Common 3D TTA includes flips, rotations, and scaling~\cite{wang2019test}. We employ a 15-view TTA: 1 base + 6 translations ($\pm$1 voxel on X/Y/Z) + 8 rotations ($\pm$2$^\circ$, $\pm$4$^\circ$ on sagittal/coronal planes), with coarse rotation alignment via normalized cross-correlation (NCC) against a template prior to TTA.

\noindent\textbf{Leakage Prevention in Medical Imaging.}
Data leakage is a pervasive issue in medical ML~\cite{esteban2020leakage,roberts2021common}. Common sources include: (1) patient-level splitting without grouping by site/scanner, (2) preprocessing statistics computed on full dataset, (3) feature selection on full data before cross-validation, (4) hyperparameter tuning on test set. We prevent leakage via: stratified group K-fold by scanner site (15 groups), per-volume z-score normalization (no dataset statistics), OOF-only calibration, and fixed hyperparameters from prior versions.